\section{Discussion}

Moving bilateral audio into a separate paper resolves a conceptual problem in the original ARC-ADHD manuscript. The core intervention no longer needs a strong theoretical defense for a speculative component, and the auditory experiment can use controls appropriate to its own mechanism.

The distinction between alternating and matched non-alternating sound is essential. Without condition B, any observed difference from quiet could be caused by rhythm, masking of environmental distraction, expectancy, novelty, general auditory orienting, or temporal pacing. The design therefore treats \(A-B\), not \(A-C\), as the central mechanistic contrast.

The relation to EMDR should remain limited. Reviews of EMDR mechanisms discuss working-memory taxation and other candidate explanations, but those findings concern different tasks, populations, and treatment procedures \parencite{landinromero2018emdr,wadji2022workingmemory}. The present study neither reproduces EMDR nor tests treatment of trauma-related memories.

A null bilateral-specific result is informative. If A and B perform similarly, the ARC-ADHD programme can retain ordinary auditory pacing if it is useful while discarding the stronger bilateral claim. If neither auditory condition helps, the component can be removed entirely.

The main future question, only after a reproducible proximal effect, would be whether the mechanism generalizes across challenge types and sessions. Clinical symptom change should not be tested prematurely.
